% Hand-drawn TikZ figure — replikace schématu Eurofound (ef25044, Figure 1)
% Replicates: https://www.eurofound.europa.eu/en/surveys-and-data/data-catalogue/...
% Edit freely — not generated.
\def\strEuPayEquitySchemaCaption{\centering Horizontální a~vertikální vliv \acp{KS} na~mzdovou rovnost na~úrovni podniku (převzato a~upraveno z~\cite{eurofound_ca_pay_equity_2026}).}%
%
\begin{figure}[htbp]
  \centering
  \resizebox{0.95\linewidth}{!}{%
  \begin{tikzpicture}[
    font=\small,
    every node/.style={align=center, inner sep=2.5pt},
    sectoral/.style       ={draw=blue!55!black, fill=blue!22, rounded corners=2.5pt, text width=2.6cm, minimum height=1.15cm, line width=0.6pt},
    sectoralNoPay/.style  ={draw=blue!55!black, fill=blue!18, rounded corners=2.5pt, text width=2.2cm, minimum height=1.05cm, line width=0.6pt},
    nosectoral/.style     ={draw=blue!55!black, fill=white,   rounded corners=2.5pt, text width=4.8cm, minimum height=1.15cm, line width=0.6pt},
    companyCA/.style      ={draw=blue!55!black, fill=blue!10, rounded corners=2.5pt, text width=1.95cm, minimum height=1.55cm, line width=0.5pt},
    company/.style        ={draw=green!45!black, fill=green!50!yellow!70, rounded corners=2.5pt, text width=2.35cm, minimum height=1.7cm, line width=0.6pt},
    bigArrow/.style       ={->, line width=2.6pt, draw=blue!50!black, >=Stealth},
    medArrow/.style       ={->, line width=1.6pt, draw=blue!50!black, >=Stealth},
    thinArrow/.style      ={->, line width=0.5pt, draw=black!75,      >=Stealth},
    dashArrow/.style      ={->, dashed, line width=0.5pt, draw=black!75, >=Stealth},
  ]
    % --- Column anchor x-coordinates (5 sloupců A--E) ---
    \def\cA{0}
    \def\cB{3.2}
    \def\cC{6.4}
    \def\cD{9.6}
    \def\cE{12.8}

    % --- Horní řádek odvětvových KS (y = 5.0) ---
    \node[sectoral]  (sA)    at (\cA,5.0) {Odvětvová \acs{KS}\\ s~mzdovou\\ strukturou};
    \node[sectoral]  (sB)    at (\cB,5.0) {Odvětvová \acs{KS}\\ s~mzdovou\\ strukturou};
    \node[nosectoral] (sNone) at ({(\cD+\cE)/2},5.0) {Bez odvětvové \acs{KS}};

    % --- Mezilehlá odvětvová KS bez mzdové struktury (nad sloupcem C, níže) ---
    \node[sectoralNoPay] (sCmid) at (\cC,3.4) {Odvětvová \acs{KS}\\ bez mzdové\\ struktury};

    % --- Podnikové KS s mzdovou strukturou (y = 2.4) ---
    \node[companyCA] (cAca) at (\cA,2.4) {Podniková\\ \acs{KS}\\ s~mzdovou\\ strukturou};
    \node[companyCA] (cDca) at (\cD,2.4) {Podniková\\ \acs{KS}\\ s~mzdovou\\ strukturou};

    % --- Spodní řádek modelových podniků A–E (y = 0) ---
    \node[company] (compA) at (\cA,0) {Podnik~A\\ dodržuje\\ vlastní \acs{KS},\\ bez výrazného\\ mzdového\\ skluzu};
    \node[company] (compB) at (\cB,0) {Podnik~B\\ následuje\\ odvětvovou\\ mzdovou strukturu,\\ s~částečným\\ mzdovým\\ skluzem};
    \node[company] (compC) at (\cC,0) {Podnik~C\\ s~vlastní mzdovou\\ strukturou,\\ příznivější\\ než odvětvové\\ minimum};
    \node[company] (compD) at (\cD,0) {Podnik~D\\ s~vlastní\\ mzdovou\\ strukturou};
    \node[company] (compE) at (\cE,0) {Podnik~E\\ bez mzdové\\ struktury};

    % --- Vertikální vliv šipkami ---
    % Sloupec A: silná šipka odvětvová KS -> podniková KS
    \draw[medArrow] (sA.south) -- (cAca.north);
    % Sloupec B: silná šipka odvětvová KS -> přímo podnik (bez podnikové KS)
    \draw[medArrow] (sB.south) -- (compB.north);
    % Sloupec C: tenká šipka odvětvová KS bez mzdové struktury -> podnik C
    \draw[thinArrow] (sCmid.south) -- (compC.north);
    % Sloupec D: čárkovaná tenká šipka odvětvová KS bez mzdové struktury -> podniková KS D
    \draw[dashArrow] (sCmid.east) -- (cDca.west);

    % --- Horní vodorovná šipka „Horizontální přístupy" ---
    \draw[bigArrow]
      ({\cA-1.9},6.45) -- ({\cE+1.9},6.45)
      node[midway, above=2pt, font=\bfseries] {Horizontální přístupy};

    % --- Levá svislá šipka „Vertikální vliv" ---
    \draw[bigArrow]
      (-2.2,5.7) -- (-2.2,-1.2)
      node[midway, left=4pt, rotate=90, anchor=south, font=\bfseries] {Vertikální vliv};
  \end{tikzpicture}}
  \caption{\strEuPayEquitySchemaCaption}
  \label{fig:eu_pay_equity_schema}
\end{figure}
